\documentclass[12pt,a4paper]{article}
\usepackage[T2A]{fontenc}
\usepackage[utf8]{inputenc}
\usepackage[russian]{babel}
\usepackage{float}
\usepackage{amsmath, amssymb}
\usepackage{geometry}
\usepackage{tikz}
\usepackage{tikz-3dplot}
\usepackage{caption}
\usetikzlibrary{arrows.meta, calc, decorations.pathmorphing}
\geometry{margin=2cm}

\newcommand{\vi}{\bar i}
\newcommand{\vj}{\bar j}
\newcommand{\vk}{\bar k}

\begin{document}

\setcounter{page}{79}

\section*{ЛЕКЦИЯ 11. ОБЗОР ПОВЕРХНОСТЕЙ ВТОРОГО ПОРЯДКА}

\textbf{Определение.} Функция $\Phi(x,y,z)$ называется алгебраическим
многочленом степени $n$, если $\Phi(x,y,z)$ есть сумма конечного
числа членов вида
$$
A\cdot x^p\cdot y^q\cdot z^r,
$$
где $p+q+r\le n$ и существует по крайней мере одно слагаемое, у
которого $p+q+r=n$.

\textbf{Определение.} Поверхность, определяемая в некоторой
декартовой системе координат алгебраическим многочленом $n$-ой
степени, называется алгебраической поверхностью $n$-ого порядка.

\subsection*{\textbf{Пример}}

$$
Ax+By+Cz+D=0
$$
--- алгебраическая поверхность 1-го порядка.

$$
\Phi(x,y,z)=a_{11}x^2+a_{22}y^2+a_{33}z^2+2a_{12}xy+
$$
$$
+2a_{13}xz+2a_{23}yz+2a_1x+2a_2y+2a_3z+a_0=0
$$
--- алгебраическая поверхность 2-го порядка (общий вид).

\subsection*{\textbf{1. Эллипсоид}}

\begin{figure}[H]
  \centering
  \begin{minipage}{0.4\textwidth}
    \centering
    \tdplotsetmaincoords{65}{115}
    \begin{tikzpicture}[scale=0.95, tdplot_main_coords, >=Latex]
      \draw[->] (0,0,0) -- (3.1,0,0) node[below left] {$x$};
      \draw[->] (0,0,0) -- (0,3.1,0) node[right] {$y$};
      \draw[->] (0,0,0) -- (0,0,3.0) node[above] {$z$};
      \node[below left] at (0,0,0) {$0$};
      \draw[thick] (0,0,0) ellipse (1.8 and 0.75);
      \draw[thick] (0,0,0) ellipse (0.75 and 1.45);
      \draw[thick, rotate around={90:(0,0,0)}] (0,0,0) ellipse (1.45 and 0.75);
      \draw[dashed] (0,0,0) -- (1.8,0,0) node[midway, below] {$a$};
      \draw[dashed] (0,0,0) -- (0,1.45,0) node[midway, right] {$b$};
      \draw[dashed] (0,0,0) -- (0,0,1.9) node[midway, left] {$c$};
    \end{tikzpicture}
    \caption*{Рисунок 11.1}
  \end{minipage}%
  \begin{minipage}{0.55\textwidth}
    $$
    \frac{x^2}{a^2}+\frac{y^2}{b^2}+\frac{z^2}{c^2}=1;
    $$
    $$
    x^2+y^2+z^2=R^2
    $$
    --- сфера.
  \end{minipage}
\end{figure}

\subsection*{\textbf{2. Однополостный гиперболоид}}

\begin{figure}[H]
  \centering
  \begin{minipage}{0.4\textwidth}
    \centering
    \tdplotsetmaincoords{65}{115}
    \begin{tikzpicture}[scale=0.9, tdplot_main_coords, >=Latex]
      \draw[->] (0,0,-2.4) -- (0,0,2.7) node[above] {$z$};
      \draw[->] (-2.7,0,0) -- (2.9,0,0) node[below left] {$x$};
      \draw[->] (0,-2.7,0) -- (0,2.9,0) node[right] {$y$};
      \draw[thick] (0,0,1.55) ellipse (1.55 and 0.55);
      \draw[thick] (0,0,0) ellipse (0.9 and 0.32);
      \draw[thick] (0,0,-1.55) ellipse (1.55 and 0.55);
      \draw[thick] (-1.55,0,1.55) .. controls (-0.85,0,0.4) and
      (-0.85,0,-0.4) .. (-1.55,0,-1.55);
      \draw[thick] (1.55,0,1.55) .. controls (0.85,0,0.4) and
      (0.85,0,-0.4) .. (1.55,0,-1.55);
      \draw[dashed] (0,0,0) -- (1.0,0,0) node[midway, below] {$a$};
      \draw[dashed] (0,0,0) -- (0,1.0,0) node[midway, right] {$b$};
    \end{tikzpicture}
    \caption*{Рисунок 11.2}
  \end{minipage}%
  \begin{minipage}{0.55\textwidth}
    $$
    \frac{x^2}{a^2}+\frac{y^2}{b^2}-\frac{z^2}{c^2}=1.
    $$
  \end{minipage}
\end{figure}

\subsection*{\textbf{3. Двуполостный гиперболоид}}

\begin{figure}[H]
  \centering
  \begin{minipage}{0.4\textwidth}
    \centering
    \tdplotsetmaincoords{65}{115}
    \begin{tikzpicture}[scale=0.9, tdplot_main_coords, >=Latex]
      \draw[->] (0,0,-2.9) -- (0,0,3.1) node[above] {$z$};
      \draw[->] (-2.5,0,0) -- (2.7,0,0) node[below left] {$x$};
      \draw[->] (0,-2.5,0) -- (0,2.7,0) node[right] {$y$};
      \draw[thick] (0,0,1.25) ellipse (0.45 and 0.18);
      \draw[thick] (0,0,2.0) ellipse (1.2 and 0.45);
      \draw[thick] (-0.45,0,1.25) .. controls (-0.75,0,1.55) .. (-1.2,0,2.0);
      \draw[thick] (0.45,0,1.25) .. controls (0.75,0,1.55) .. (1.2,0,2.0);
      \draw[thick] (0,0,-1.25) ellipse (0.45 and 0.18);
      \draw[thick] (0,0,-2.0) ellipse (1.2 and 0.45);
      \draw[thick] (-0.45,0,-1.25) .. controls (-0.75,0,-1.55) .. (-1.2,0,-2.0);
      \draw[thick] (0.45,0,-1.25) .. controls (0.75,0,-1.55) .. (1.2,0,-2.0);
      \node[right] at (0,0,1.25) {$c$};
      \node[right] at (0,0,-1.25) {$-c$};
    \end{tikzpicture}
    \caption*{Рисунок 11.3}
  \end{minipage}%
  \begin{minipage}{0.55\textwidth}
    $$
    \frac{x^2}{a^2}+\frac{y^2}{b^2}-\frac{z^2}{c^2}=-1.
    $$
  \end{minipage}
\end{figure}

\subsection*{\textbf{4. Конус}}

\begin{figure}[H]
  \centering
  \begin{minipage}{0.4\textwidth}
    \centering
    \tdplotsetmaincoords{65}{115}
    \begin{tikzpicture}[scale=0.9, tdplot_main_coords, >=Latex]
      \draw[->] (0,0,-2.4) -- (0,0,2.7) node[above] {$z$};
      \draw[->] (-2.6,0,0) -- (2.8,0,0) node[below left] {$x$};
      \draw[->] (0,-2.6,0) -- (0,2.8,0) node[right] {$y$};
      \draw[thick] (0,0,1.8) ellipse (1.45 and 0.55);
      \draw[thick] (0,0,-1.8) ellipse (1.45 and 0.55);
      \draw[thick] (-1.45,0,1.8) -- (0,0,0) -- (1.45,0,1.8);
      \draw[thick] (-1.45,0,-1.8) -- (0,0,0) -- (1.45,0,-1.8);
      \node[below left] at (0,0,0) {$0$};
    \end{tikzpicture}
    \caption*{Рисунок 11.4}
  \end{minipage}%
  \begin{minipage}{0.55\textwidth}
    $$
    \frac{x^2}{a^2}+\frac{y^2}{b^2}-\frac{z^2}{c^2}=0.
    $$
  \end{minipage}
\end{figure}

\subsection*{\textbf{5. Эллиптический параболоид}}

\begin{figure}[H]
  \centering
  \begin{minipage}{0.4\textwidth}
    \centering
    \tdplotsetmaincoords{65}{115}
    \begin{tikzpicture}[scale=0.9, tdplot_main_coords, >=Latex]
      \draw[->] (0,0,0) -- (0,0,3.2) node[above] {$z$};
      \draw[->] (-2.5,0,0) -- (2.7,0,0) node[below left] {$x$};
      \draw[->] (0,-2.5,0) -- (0,2.7,0) node[right] {$y$};
      \draw[thick] (0,0,0.7) ellipse (0.7 and 0.25);
      \draw[thick] (0,0,1.45) ellipse (1.2 and 0.45);
      \draw[thick] (0,0,2.35) ellipse (1.7 and 0.65);
      \draw[thick] (-0.7,0,0.7) .. controls (-1.1,0,1.3) .. (-1.7,0,2.35);
      \draw[thick] (0.7,0,0.7) .. controls (1.1,0,1.3) .. (1.7,0,2.35);
      \node[below left] at (0,0,0) {$0$};
    \end{tikzpicture}
    \caption*{Рисунок 11.5}
  \end{minipage}%
  \begin{minipage}{0.55\textwidth}
    $$
    \frac{x^2}{p}+\frac{y^2}{q}=2z,
    $$
    $$
    p>0,
    $$
    $$
    q>0.
    $$
  \end{minipage}
\end{figure}

\subsection*{\textbf{6. Гиперболический параболоид}}

\begin{figure}[H]
  \centering
  \begin{minipage}{0.4\textwidth}
    \centering
    \tdplotsetmaincoords{65}{115}
    \begin{tikzpicture}[scale=0.85, tdplot_main_coords, >=Latex]
      \draw[->] (-2.7,0,0) -- (2.9,0,0) node[below left] {$x$};
      \draw[->] (0,-2.7,0) -- (0,2.9,0) node[right] {$y$};
      \draw[->] (0,0,-1.2) -- (0,0,2.8) node[above] {$z$};
      \draw[thick] (-2.3,0,1.0) .. controls (-1.0,0,0.05) and
      (1.0,0,0.05) .. (2.3,0,1.0);
      \draw[thick] (0,-2.3,-0.9) .. controls (0,-1.0,-0.05) and
      (0,1.0,-0.05) .. (0,2.3,-0.9);
      \draw[thick] (-2.2,-1.7,0.45) -- (2.2,-1.7,1.15);
      \draw[thick] (-2.2,1.7,1.15) -- (2.2,1.7,0.45);
      \draw[dashed] (-2.2,0,0) -- (2.2,0,0);
      \draw[dashed] (0,-2.2,0) -- (0,2.2,0);
      \node[below left] at (0,0,0) {$0$};
    \end{tikzpicture}
    \caption*{Рисунок 11.6}
  \end{minipage}%
  \begin{minipage}{0.55\textwidth}
    $$
    \frac{x^2}{p}-\frac{y^2}{q}=2z;
    $$
    $$
    p>0;
    $$
    $$
    q>0.
    $$
  \end{minipage}
\end{figure}

\subsection*{\textbf{7. Эллиптический цилиндр}}

\begin{figure}[H]
  \centering
  \begin{minipage}{0.4\textwidth}
    \centering
    \tdplotsetmaincoords{65}{115}
    \begin{tikzpicture}[scale=0.9, tdplot_main_coords, >=Latex]
      \draw[->] (-2.4,0,0) -- (2.6,0,0) node[below left] {$x$};
      \draw[->] (0,-2.4,0) -- (0,2.6,0) node[right] {$y$};
      \draw[->] (0,0,-0.5) -- (0,0,3.5) node[above] {$z$};
      \draw[thick] (0,0,3) ellipse (1.25 and 0.55);
      \draw[thick] (0,0,0) ellipse (1.25 and 0.55);
      \draw[thick] (1.25,0,0) -- (1.25,0,3);
      \draw[thick] (-1.25,0,0) -- (-1.25,0,3);
      \draw[dashed] (0,0,0) -- (1.25,0,0) node[midway, below] {$a$};
      \draw[dashed] (0,0,0) -- (0,0.75,0) node[midway, right] {$b$};
      \draw[dashed] (0,0,0) -- (0,0,3);
    \end{tikzpicture}
    \caption*{Рисунок 11.7}
  \end{minipage}%
  \begin{minipage}{0.55\textwidth}
    $$
    \frac{x^2}{a^2}+\frac{y^2}{b^2}=1.
    $$
  \end{minipage}
\end{figure}

\subsection*{\textbf{8. Гиперболический цилиндр}}

\begin{figure}[H]
  \centering
  \begin{minipage}{0.4\textwidth}
    \centering
    \tdplotsetmaincoords{65}{115}
    \begin{tikzpicture}[scale=0.82, tdplot_main_coords, >=Latex]
      \draw[->] (-2.9,0,0) -- (3.1,0,0) node[below left] {$x$};
      \draw[->] (0,-2.8,0) -- (0,3.0,0) node[right] {$y$};
      \draw[->] (0,0,-0.5) -- (0,0,3.4) node[above] {$z$};
      \draw[thick] plot[smooth] coordinates {(-2.2,0,0) (-1.5,0.35,0)
      (-1.2,0.9,0)};
      \draw[thick] plot[smooth] coordinates {(-2.2,0,3) (-1.5,0.35,3)
      (-1.2,0.9,3)};
      \draw[thick] plot[smooth] coordinates {(2.2,0,0) (1.5,-0.35,0)
      (1.2,-0.9,0)};
      \draw[thick] plot[smooth] coordinates {(2.2,0,3) (1.5,-0.35,3)
      (1.2,-0.9,3)};
      \draw[thick] (-2.2,0,0) -- (-2.2,0,3);
      \draw[thick] (-1.2,0.9,0) -- (-1.2,0.9,3);
      \draw[thick] (2.2,0,0) -- (2.2,0,3);
      \draw[thick] (1.2,-0.9,0) -- (1.2,-0.9,3);
      \draw[dashed] (-1.2,-0.9,0) -- (-1.2,-0.9,3);
      \draw[dashed] (1.2,0.9,0) -- (1.2,0.9,3);
      \draw[dashed] (0,-2.0,0) -- (0,2.0,0);
    \end{tikzpicture}
    \caption*{Рисунок 11.8}
  \end{minipage}%
  \begin{minipage}{0.55\textwidth}
    $$
    \frac{x^2}{a^2}-\frac{y^2}{b^2}=1.
    $$
  \end{minipage}
\end{figure}

\subsection*{\textbf{9. Параболический цилиндр}}

\begin{figure}[H]
  \centering
  \begin{minipage}{0.4\textwidth}
    \centering
    \tdplotsetmaincoords{65}{115}
    \begin{tikzpicture}[scale=0.85, tdplot_main_coords, >=Latex]
      \draw[->] (-2.6,0,0) -- (2.8,0,0) node[below left] {$x$};
      \draw[->] (0,-2.4,0) -- (0,2.8,0) node[right] {$y$};
      \draw[->] (0,0,-0.4) -- (0,0,3.4) node[above] {$z$};
      \draw[thick] plot[smooth] coordinates {(-1.9,-1.1,0)
      (-0.7,-0.4,0) (0,0,0) (-0.7,0.4,0) (-1.9,1.1,0)};
      \draw[thick] plot[smooth] coordinates {(-1.9,-1.1,3)
      (-0.7,-0.4,3) (0,0,3) (-0.7,0.4,3) (-1.9,1.1,3)};
      \draw[thick] (-1.9,-1.1,0) -- (-1.9,-1.1,3);
      \draw[thick] (-1.9,1.1,0) -- (-1.9,1.1,3);
      \draw[dashed] (0,0,0) -- (0,0,3);
      \draw[dashed] (-1.9,0,0) -- (0,0,0);
    \end{tikzpicture}
    \caption*{Рисунок 11.9}
  \end{minipage}%
  \begin{minipage}{0.55\textwidth}
    $$
    y^2=2px.
    $$
  \end{minipage}
\end{figure}

\textbf{Определение.} Прямолинейной образующей поверхности называется
прямая линия, целиком лежащая на данной поверхности.

\textbf{Замечание.} Поверхности под номерами 1, 3, 5 прямолинейных
образующих не имеют, а другие имеют.

\end{document}
